% ==================================================================
%  Implementation.tex  --  Section 4 (implementation and evaluation)
%  Cross-references used here:
%    \label{sec:dbscan} \label{sec:post}   (Sections 2-3)
%  Table tab:results: runs of 2026-11-13/14 with alpha_l = ceil(log2(1.1 n))
%  (logs desilo/{hepta,tetra,lsun,twodiamonds,longsquare,s1,wbc,swiss}.txt,
%  debug_fhe_vs_np_*.csv).  Time = Normalize + Core + Label_Propagation
%  ("합계") as printed by the test script, in seconds: hepta 609.06,
%  tetra 965.93, lsun 1258.26, twodiamonds 1586.09, longsquare 1865.99,
%  s1 8180.36, banknote 808.54, wbc 3791.58 (server only).
% ==================================================================
\section{Implementation and Evaluation}
\label{sec:eval}

\subsection{Setup}
\label{sec:setup}

We implemented the protocol in Python with DesiloFHE~1.9.0, with the
server side running on an NVIDIA GeForce RTX 4090 GPU with 24\,GB of
memory, and evaluated it on the datasets of Table~\ref{tab:setup}: six
synthetic benchmarks and two real datasets from the application areas of
Section~\ref{sec:intro}, outlier detection on financial data and the
analysis of healthcare data.

\begin{table}[t]
\centering
\caption{Datasets.}
\label{tab:setup}
\renewcommand{\arraystretch}{1.1}
\begin{tabular}{@{}l@{\hspace{1.5em}}r@{\hspace{1.5em}}r@{\hspace{2em}}p{0.55\linewidth}@{}}
\toprule
Dataset & $n$ & $d$ & Purpose \tabularnewline
\midrule
hepta        &  212 & 3 & \raggedright baseline: well-separated convex
  clusters \tabularnewline
tetra        &  400 & 3 & \raggedright few core points: most points are
  border points \tabularnewline
lsun         &  400 & 2 & \raggedright two bars forming an L and a round
  cluster, with unequal spreads and gaps \tabularnewline
twodiamonds  &  800 & 2 & \raggedright two touching clusters, separated
  only by density \tabularnewline
longsquare   &  900 & 2 & \raggedright clusters of mixed shapes: squares
  and long, thin bars \tabularnewline
s1           & 5000 & 2 & \raggedright largest $n$; 15 Gaussian clusters,
  some of them nearly touching \tabularnewline
\midrule
banknote     &  200 & 6 & \raggedright finance: genuine and counterfeit notes
  \tabularnewline
wbc          &  683 & 9 & \raggedright healthcare: benign and malignant tumors
  \tabularnewline
\bottomrule
\end{tabular}
\end{table}

banknote~\cite{flury1988multivariate} holds six measurements, in
millimeters, of $100$ genuine and $100$ counterfeit Swiss banknotes. wbc
is the Wisconsin Breast Cancer (Original)
dataset~\cite{wolberg1990breast}, in which nine cell features, each
graded from $1$ to $10$, describe benign and malignant breast tumors. We
remove its $16$ samples with a missing value, which leaves $683$.

For banknote, the client first subtracts from each feature its minimum.
The raw values range from $7.2$ to $216.3$ although no feature varies by
more than $5.5$, so the single scale $\sigma$ of
Section~\ref{sec:client} would be set by this common offset and
$\varepsilon_n^2$ would fall below the margin $\mu$ of
Eq.~\eqref{eq:theta}. A translation leaves every pairwise distance, and
hence the clustering, unchanged.
% Input file: dataset/Other_cluster/swiss_banknote_shifted.arff (rows of
%   the R package mclust, data(banknote), minus the column minima 213.8,
%   129, 129, 7.2, 7.7, 137.8).
% wbc: dataset/Other_cluster/wbc_original.arff, UCI id 15, 699 rows; the
%   16 removed rows have '?' in Bare Nuclei; ID column removed; no scaling.

\subsection{Clustering Accuracy}
\label{sec:accuracy}

On all eight datasets the encrypted run returns the same clustering as
plaintext DBSCAN at the same $\varepsilon$ and $\mathit{MinPts}$ that
assigns border points by the maximum rule of Section~\ref{sec:dbscan}:
the same clusters and the same noise points.
Table~\ref{tab:results} reports the server time and the adjusted Rand
index (ARI) of the encrypted result, with noise as one group, against the
ground truth and against plaintext DBSCAN.
% Plaintext reference: sklearn.cluster.DBSCAN (scikit-learn 1.9.1),
%   DBSCAN(eps = eps_n, min_samples = MinPts) on the normalized points
%   (test_Server_main_call_dataset.py, Step 5).  min_samples counts the
%   point itself and the radius is closed (<= eps), as in Section 2.
%   Its labels equal those on the raw data (uniform scaling).

The neighbor counts behave as Eq.~\eqref{eq:sandwich} predicts. In every
run, the encrypted count of each point lies between the plaintext counts
at $\varepsilon_n$ and $\varepsilon_{\mathrm{hi}}$, and the core points
coincide with those of plaintext DBSCAN except for one point of
longsquare. That point has $\mathit{MinPts} - 1$ neighbors within
$\varepsilon_n$ and one more at a distance between $\varepsilon_n$ and
$\varepsilon_{\mathrm{hi}}$, so the encrypted run counts
$\mathit{MinPts}$ neighbors and makes it a core point, as
Theorem~\ref{thm:sandwich} allows. The point is a border point of the
same cluster in plaintext DBSCAN, so the clustering does not change.

\begin{table}[t]
\centering
\caption{Server running time and ARI of the encrypted result against the
  ground-truth labels and against plaintext DBSCAN.}
\label{tab:results}
\begin{tabular}{@{}lrrr@{}}
\toprule
 & & \multicolumn{2}{c@{}}{ARI (\%)} \\
\cmidrule(l){3-4}
Dataset & Time (min) & vs.\ ground truth & vs.\ plaintext \\
\midrule
hepta        & $ 10.2$ & $100.00$ & $100.00$ \\
tetra        & $ 16.1$ & $100.00$ & $100.00$ \\
lsun         & $ 21.0$ & $100.00$ & $100.00$ \\
twodiamonds  & $ 26.4$ & $ 99.75$ & $100.00$ \\
longsquare   & $ 31.1$ & $ 95.02$ & $100.00$ \\
s1           & $136.3$ & $ 97.89$ & $ 99.96$ \\
\midrule
banknote     & $ 13.5$ & $ 94.06$ & $100.00$ \\
wbc          & $ 63.2$ & $ 87.98$ & $100.00$ \\
\bottomrule
\end{tabular}

\vspace{2pt}
\parbox{\linewidth}{\footnotesize Plaintext: DBSCAN of
  scikit-learn~\cite{pedregosa2011scikit}, same $\varepsilon_n$ and
  $\mathit{MinPts}$.}
\end{table}

Where the ARI against the ground truth falls below $100\%$, the cause
lies in DBSCAN itself and not in the encryption. DBSCAN reports some
points as noise that the ground truth assigns to a cluster, $1$ point of
twodiamonds, $35$ of longsquare and $85$ of s1, and on s1 it places two
border points in the cluster of a neighboring class; plaintext DBSCAN
with the maximum rule gives exactly the same result.
Against plaintext DBSCAN, the ARI is $100\%$ except on s1. The noise points
coincide on all eight datasets, but scikit-learn gives a border point to
the cluster that reaches it first (Section~\ref{sec:dbscan}), and of the
seven border points of s1 that neighbor two clusters, it places one
differently from the maximum rule. The encrypted run puts it in its true
class and scores $97.89\%$ against the ground truth, scikit-learn
$97.84\%$.

\paragraph{\emph{Real data}}
The two real datasets show the two properties for which
Section~\ref{sec:intro} prefers DBSCAN to K-means. On banknote, DBSCAN
finds the two kinds of notes without being given the number of clusters:
one cluster holds $97$ genuine notes, the other $98$ counterfeit notes
and one genuine note, and four notes, two of each kind, are reported as
noise. On wbc, the noise label itself is the result. DBSCAN finds a
single cluster, which holds $431$ of the $444$ benign samples, and
reports $231$ of the $239$ malignant samples as noise, that is, as
outliers of the benign population. In both cases the server works on
ciphertexts alone, so such financial or medical records can be analyzed
without being revealed.

\paragraph{\emph{Running time}}
Every run except s1 and wbc finishes in about half an hour or less, and
s1, with $5000$ points, in $2.3$ hours. The time follows the number of
evaluations of $\mathsf{Sgn}$ rather than $n$ alone. Neighbor Check
evaluates $\mathsf{Sgn}$ once per offset, $k_{\max}$ times in total, and
takes between $14\%$ and $43\%$ of the time. Label Propagation calls
$\mathsf{Max}$ at most $2g \lceil \log_2(m+1) \rceil$ times
(Section~\ref{sec:lp}) in each of the $R$ rounds and in the final pass
to the non-core points, and takes almost all of the rest.
% Timing split from the logs (Normalize / Core / Label_Propagation, s):
%   hepta 229.37/1.85/377.84, tetra 414.37/1.88/549.68,
%   lsun 313.82/1.89/942.55, twodiamonds 215.70/1.96/1368.43,
%   longsquare 273.90/1.94/1590.15, s1 1868.97/2.44/6308.95,
%   banknote 241.83/1.83/564.88, wbc 1346.76/1.93/2442.89.
% Real data, confusion of the encrypted run (2026-11-13 logs):
%   banknote: cluster A 97 genuine; cluster B 98 counterfeit + 1 genuine;
%   noise 2 genuine + 2 counterfeit.  wbc: cluster 431 benign + 8
%   malignant; noise 13 benign + 231 malignant.
% Checked on the 2026-11-13/14 runs (all eight logs):
%   the partition equals plaintext max-rule DBSCAN (labels = largest core
%   slot + 1) on all eight datasets, noise sets equal scikit-learn's, and
%   the core sets equal plaintext DBSCAN's except longsquare index 399.
%   banknote: 40 points count one neighbor more than plaintext (20 pairs
%   at a distance between eps and eps_hi); wbc: all counts equal (no pair
%   between eps and eps_hi).  On s1,
%   scikit-learn assigns one of the seven ambiguous border points
%   differently (index 1660; the encrypted run puts it in its true class),
%   hence ARI(FHE, scikit-learn) = 99.96.  On twodiamonds one border point
%   (index 120) neighbors both clusters; both assign it to its true class.
%   s1 misassigned border points of the encrypted run: 208 (ambiguous,
%   true 13 -> cluster of class 14) and 3561 (core neighbors only in the
%   cluster of class 8, true 9).
%   Sandwich of Eq. (sandwich) holds for every point with eps_hi from the
%   logged alpha_d; extra core point on longsquare: index 399.


\subsection{Comparison with Prior Work}
\label{sec:compare}

Table~\ref{tab:compare} places our results next to those reported for the
schemes of Section~\ref{sec:related} on the datasets they share with ours.
The values of the other schemes are taken from the original papers, each
obtained with the parameters chosen there, so the table compares clustering
quality against the ground truth rather than identical runs. The HE schemes
for mean-shift~\cite{cheon2019towards} and K-means~\cite{yang2024optimized}
report accuracy, and the MPC-based DBSCAN
schemes~\cite{bozdemir2021privacy,fu2024ppa} report the ARI; we give our
result in the same metric. For our scheme, accuracy is the fraction of
points whose cluster matches the ground-truth class under the best
one-to-one matching, with noise counted as an error.

\begin{table}[t]
\centering
\caption{Comparison with prior privacy-preserving clustering on shared
  datasets. Values of the other schemes are as reported in the original
  papers, each obtained with the parameters chosen there; a dash marks a
  value that is not reported. On lsun, the mean-shift scheme returns four
  clusters instead of three; their higher silhouette coefficient reflects
  more compact clusters rather than agreement with the ground truth, which
  plaintext mean-shift does not recover either
  (Section~\ref{sec:compare}).}
\label{tab:compare}
\begin{tabular}{lccc}
\toprule
\multicolumn{4}{l}{(a) Accuracy (\%), HE-based clustering} \\
\midrule
Dataset & Cheon et al.~\cite{cheon2019towards}
  & Yang et al.~\cite{yang2024optimized} & Ours \\
\midrule
hepta        & 100  & 80.2  & 100 \\
lsun         & --   & 71.2  & 100 \\
tetra        & 100  & 96.8  & 100 \\
twodiamonds  & 99.0 & 100   & 99.9 \\
\midrule
\multicolumn{4}{l}{(b) ARI (\%), MPC-based DBSCAN} \\
\midrule
Dataset & Bozdemir et al.~\cite{bozdemir2021privacy}
  & Fu et al.~\cite{fu2024ppa} & Ours \\
\midrule
lsun         & 100  & 100   & 100 \\
s1           & 97.6 & 97.8  & 97.9 \\
longsquare   & --   & 93.4  & 95.0 \\
\bottomrule
\end{tabular}
\end{table}
% Sources of the reported values:
%   Cheon et al.  ePrint 2019/465, Table 2 (hepta 212/212, tetra 400/400,
%                 twodiamonds 792/800; lsun not reported, 4 clusters found).
%   Yang et al.   ePrint 2024/1141, Table 2, column "Ours" (encrypted).
%   Bozdemir et al. ASIACCS'21, Table 1 (lsun 1.0, s1 0.9757).
%   Fu et al.     TDSC'24, Table VIII (lsun, longsquare, rho = 0.1) and
%                 Table XI (s1, rho = 0.1).
%   Ours (2026-11 runs): accuracy hepta 100, lsun 100, tetra 100,
%   twodiamonds 99.875 (799/800); ARI lsun 100, s1 97.886, longsquare 95.020.

% Table 5 (Section 5, tab:leakage). Defined here, not in SecurityAnalysis.tex:
% a two-column float can only appear on a page after the one where it is
% defined, so defining it here places it on the page where Section 5 starts.
\begin{table*}[t]
\centering
\caption{Values the server receives in the clear.}
\label{tab:leakage}
\small
\renewcommand{\arraystretch}{1.15}
% \raggedright in each p-cell avoids stretched spacing in the narrow
% columns; rows therefore end with \tabularnewline instead of \\.
\begin{tabular}{@{}l p{0.16\linewidth} p{0.25\linewidth} p{0.35\linewidth}@{}}
\toprule
Value & Description & Role at the server & What it reveals about the data
  \tabularnewline
\midrule
$n$ & \raggedright number of points
    & \raggedright slot layout, loop bounds, precisions $\alpha$
    & \raggedright size of the dataset \tabularnewline
$d$ & \raggedright dimension
    & \raggedright number of input ciphertexts
    & \raggedright dimension of the dataset \tabularnewline
$\epsn$, $\MinPts$
    & \raggedright DBSCAN parameters
    & \raggedright constants in Eqs.~(\ref{eq:cmpinput}) and~(\ref{eq:coreinput})
    & \raggedright $\varepsilon$ relative to the normalization scale;
      density threshold \tabularnewline
$\kmax$
    & \raggedright operation range
    & \raggedright offsets in Algorithms~\ref{alg:neighbor}--\ref{alg:lp}
    & \raggedright one less than the largest number of points in a window
      of width $\epsn$ along the first principal component, capped at
      $\lfloor n/2 \rfloor$ \tabularnewline
$R$ & \raggedright number of rounds
    & \raggedright rounds in Algorithm~\ref{alg:lp}
    & \raggedright a value chosen by the client; with the grid estimate,
      the largest geodesic diameter of a cell group in units of $\epsn$
      \tabularnewline
\bottomrule
\end{tabular}
\end{table*}

Against the HE schemes, the differences on hepta and lsun stem from the
algorithms rather than from the encryption. There only K-means falls
behind, and Yang et al.~\cite{yang2024optimized} obtain nearly the same
accuracy with plaintext K-means on these datasets. Mean-shift reaches
$100\%$ on hepta; lsun is discussed below. Against the MPC-based schemes, our
ARI is on par. Bozdemir et al.~\cite{bozdemir2021privacy} evaluate exact
DBSCAN, so the small difference on s1 reflects a different choice of
parameters rather than a loss of accuracy on either side. PPA-DBSCAN
returns an approximation of DBSCAN: on longsquare its ARI lies below the
94.9 that Fu et al.~\cite{fu2024ppa} report for exact DBSCAN, whereas our
scheme returns the clustering of plaintext DBSCAN.

On lsun, Cheon et al.~\cite{cheon2019towards} obtain four clusters
instead of the three of the ground truth and therefore report no
accuracy. They note instead that their four clusters have a higher
silhouette coefficient than the ground-truth labels, $0.577$ against
$0.443$. The silhouette coefficient, however, measures how compact and
well separated the clusters are, not whether they match the ground truth.
lsun consists of two elongated bars that form an L and a round cluster,
and dividing an elongated cluster into more compact parts raises the
coefficient even as the partition moves away from the ground truth. The
four clusters are a property of mean-shift itself rather than of its
encrypted evaluation: mean-shift assigns each point to a mode of the
density, and a bar of nearly uniform density need not have a single mode.
In our plaintext runs of the mean-shift implementation of
scikit-learn~\cite{pedregosa2011scikit}, no bandwidth between $0.2$ and
$3.0$ recovered the ground truth, the best ARI being $0.65$. DBSCAN
instead links points through chains of neighbors, so it follows each bar
along its length, and both plaintext DBSCAN and our encrypted run recover
the three clusters of lsun.
% Plaintext mean-shift check (2026-11-11): sklearn.cluster.MeanShift on
%   dataset/Other_cluster/lsun.arff, bandwidth 0.20:0.05:3.00; best ARI
%   0.648 (bw 1.00-1.10, 5 clusters); ARI = 1 never reached.